% =====================================================================
% Warm Folio - preámbulo LaTeX del libro (salida PDF, LuaLaTeX).
% Valores tomados de styles/design/tokens.json.
% Se incluye con format.pdf.include-in-header.
%
% Nota de orden: Quarto carga amsthm, caption, float y los entornos de
% teorema DESPUÉS de este archivo. Todo lo que depende de ellos va dentro
% de \AtBeginDocument, con \makeatletter FUERA del gancho (los tokens del
% argumento se leen aquí, no al ejecutarse) y con # duplicado.
% =====================================================================

\usepackage{iftex}
\usepackage{xcolor}
\usepackage{etoolbox}
\usepackage{tcolorbox}
\tcbuselibrary{skins,breakable}
\usepackage{setspace}
\usepackage{marginnote}
\usepackage{scrlayer-scrpage}
\usepackage{ragged2e}
\usepackage{tikz}
% ---------- Colores (tokens) ----------
\definecolor{ink}{HTML}{2A231D}
\definecolor{muted}{HTML}{6B5F52}
\definecolor{paper}{HTML}{FDFBF7}
\definecolor{tint}{HTML}{F7EDE5}
\definecolor{tintstrong}{HTML}{F0E0D6}
\definecolor{codebg}{HTML}{FBF6F1}
\definecolor{rule}{HTML}{DED0C4}
\definecolor{rulesoft}{HTML}{E7DAD0}
\definecolor{accent}{HTML}{B0472A}
\definecolor{accentdark}{HTML}{8E3620}
\definecolor{olive}{HTML}{4E6E38}
\definecolor{steel}{HTML}{3D647F}
\definecolor{shadecolor}{HTML}{FBF6F1}

\color{ink}
\setlength{\emergencystretch}{2em}
\frenchspacing
\widowpenalty=10000
\clubpenalty=10000

% ---------- Títulos ----------
\KOMAoptions{chapterprefix=true,headings=normal}
\setkomafont{chapterprefix}{\sffamily\bfseries\fontsize{9}{11}\selectfont\color{accent}}
\renewcommand*{\chapterformat}{\MakeUppercase{\chapapp\nobreakspace\thechapter}}
\setkomafont{chapter}{\rmfamily\bfseries\fontsize{26}{29}\selectfont\color{ink}}
\setkomafont{section}{\sffamily\bfseries\fontsize{14}{17}\selectfont\color{ink}}
\setkomafont{subsection}{\sffamily\bfseries\fontsize{11.5}{14}\selectfont\color{ink}}
\setkomafont{subsubsection}{\sffamily\bfseries\fontsize{10.5}{13}\selectfont\color{muted}}

% ---------- Encabezados y folios ----------
\clearpairofpagestyles
\automark[section]{chapter}
\setkomafont{pagehead}{\sffamily\fontsize{9}{11}\selectfont\color{muted}}
\setkomafont{pagenumber}{\sffamily\fontsize{9}{11}\selectfont\color{muted}}
\newcommand{\wfshorttitle}{Macroeconometría aplicada con Python}
\lehead{\pagemark\quad\MakeUppercase{\wfshorttitle}}
\rohead{\MakeUppercase{\rightmark}\quad\pagemark}
\cfoot{}
\renewcommand*{\chapterpagestyle}{plain}

% Con \flushbottom (el estilo por omisión de scrbook a dos caras) LaTeX estira
% los espacios verticales cuando una figura flotante se pospone, y las páginas
% con mucho código quedan con huecos enormes entre bloques. \raggedbottom deja
% el sobrante al pie, que es lo habitual en un libro técnico.
\raggedbottom

% ---------- Bloques de código ----------
% fvextra parte las líneas largas de los bloques de código. Quarto no lo carga
% por su cuenta en este proyecto, así que se carga aquí (después de fancyvrb,
% que ya está en el preámbulo) y se fija el comportamiento de corte:
% breakanywhere evita que un identificador largo se salga del bloque de texto.
\usepackage{fvextra}
\fvset{breaklines=true,breakanywhere=true,breakindent=1.2em,
  breaksymbolleft={\color{muted}\footnotesize\ensuremath{\hookrightarrow}},
  breaksymbolindentleft=0pt,breaksymbolsepleft=0.4em}

% El entorno Shaded de pandoc usa framed/snugshade, que al partirse entre
% páginas deja una franja teñida vacía. Se redefine con tcolorbox partible,
% que corta limpio y conserva el mismo fondo sin borde.
\AtBeginDocument{%
  \renewenvironment{Shaded}{%
    \begin{tcolorbox}[breakable, enhanced jigsaw, colback=codebg,
      frame hidden, boxrule=0pt, sharp corners, boxsep=0pt,
      left=6pt, right=6pt, top=5pt, bottom=5pt,
      grow to left by=6pt, grow to right by=6pt,
      before skip=8pt, after skip=8pt]}{\end{tcolorbox}}}

% ---------- Notas al margen ----------
\renewcommand*{\marginfont}{\sffamily\fontsize{8}{10.5}\selectfont\color{muted}\RaggedRight}

% ---------- Etiquetas de los bloques ----------
\newcommand{\wfheadcolor}{ink}
% \protected: newtheoremstyle hace xdef y así el nombre no se expande allí.
\protected\def\wfheadfont{\sffamily\bfseries\fontsize{8.5}{10.5}\selectfont\color{\wfheadcolor}}
\newcommand{\wfsetheadcolor}[2]{\AtBeginEnvironment{#1}{\renewcommand{\wfheadcolor}{#2}}}
\newcommand{\wflabel}[2]{\noindent{\sffamily\bfseries\fontsize{8.5}{10.5}\selectfont\color{#1}#2}\par\vspace{3pt}}

\makeatletter
\newcommand{\wfblockstyle}{%
  \newtheoremstyle{wfblock}%
    {8pt}{8pt}%            espacio arriba / abajo
    {\normalfont}%         cuerpo (sin cursiva)
    {0pt}%                 sangría
    {\wfheadfont}%         tipo de la etiqueta
    {}%                    sin puntuación tras la etiqueta
    {\newline}%            la etiqueta ocupa su propia línea
    {\thmname{\MakeUppercase{##1}}\thmnumber{\ ##2}%
     \thmnote{\ \textperiodcentered\ \MakeUppercase{##3}}}%
  \let\th@plain\th@wfblock
  \let\th@definition\th@wfblock
  \let\th@remark\th@wfblock}
\makeatother

% ---------- Cajas de los bloques ----------
% Estilo "sólo filetes" de los bloques de resultado, como clave de tcolorbox
% (un \newcommand no sirve: pgfkeys no lo expande antes de leer las claves).
\tcbset{wfsolofiletes/.style={
  blanker, breakable, boxsep=0pt,
  left=0pt, right=0pt, top=7pt, bottom=7pt,
  borderline north={1.2pt}{0pt}{ink},
  borderline south={0.4pt}{0pt}{ink},
  before skip=12pt, after skip=12pt}}

\AtBeginDocument{%
  \wfblockstyle
  % color de la etiqueta por bloque
  \wfsetheadcolor{conjecture}{accentdark}%   Supuesto
  \wfsetheadcolor{definition}{ink}%          Definición
  \wfsetheadcolor{theorem}{ink}%             Resultado
  \wfsetheadcolor{proposition}{ink}%
  \wfsetheadcolor{lemma}{ink}%
  \wfsetheadcolor{corollary}{ink}%
  \wfsetheadcolor{example}{steel}%           Ejemplo / aplicación
  \wfsetheadcolor{exercise}{muted}%          Ejercicio
  % Supuesto: el bloque más fuerte del sistema.
  \tcolorboxenvironment{conjecture}{
    blanker, breakable, colback=tintstrong, boxsep=0pt,
    left=10pt, right=10pt, top=8pt, bottom=8pt,
    borderline west={2.5pt}{0pt}{accent},
    before skip=12pt, after skip=12pt}
  % Definición: fondo tenue.
  \tcolorboxenvironment{definition}{
    blanker, breakable, colback=tint, boxsep=0pt,
    left=10pt, right=10pt, top=8pt, bottom=8pt,
    before skip=12pt, after skip=12pt}
  % Resultado: sólo filetes.
  \tcolorboxenvironment{theorem}{wfsolofiletes}
  \tcolorboxenvironment{proposition}{wfsolofiletes}
  \tcolorboxenvironment{lemma}{wfsolofiletes}
  \tcolorboxenvironment{corollary}{wfsolofiletes}
  % Ejemplo: caja blanca con filete.
  \tcolorboxenvironment{example}{
    blanker, breakable, colback=white, boxsep=0pt,
    left=10pt, right=10pt, top=8pt, bottom=8pt,
    borderline={0.4pt}{0pt}{rule},
    before skip=12pt, after skip=12pt}
  % Ejercicio: filetes arriba y abajo.
  \tcolorboxenvironment{exercise}{
    blanker, breakable, boxsep=0pt,
    left=0pt, right=0pt, top=6pt, bottom=6pt,
    borderline north={0.4pt}{0pt}{rule},
    borderline south={0.4pt}{0pt}{rule},
    before skip=10pt, after skip=10pt}}

% ---------- Bloques sin numeración (los produce filters/blocks.lua) ----------
\newtcolorbox{intuicionblock}{
  blanker, breakable, boxsep=0pt,
  left=10pt, right=0pt, top=6pt, bottom=6pt,
  borderline west={1.5pt}{0pt}{olive,dashed},
  before skip=12pt, after skip=12pt,
  before upper={\wflabel{olive}{INTUICIÓN}}}

\newtcolorbox{cuidadoblock}{
  blanker, breakable, boxsep=0pt,
  left=10pt, right=0pt, top=6pt, bottom=6pt,
  borderline north={0.4pt}{0pt}{accentdark},
  borderline west={2.5pt}{0pt}{accentdark},
  before skip=12pt, after skip=12pt,
  before upper={\wflabel{accentdark}{CUIDADO}}}

\newtcolorbox{solucionblock}{
  blanker, breakable, boxsep=0pt,
  left=0pt, right=0pt, top=2pt, bottom=6pt,
  before skip=10pt, after skip=12pt,
  fontupper=\small,
  before upper={\wflabel{muted}{SOLUCIÓN}}}

\newtcolorbox{preguntablock}{
  blanker, breakable, boxsep=0pt,
  left=12pt, right=0pt, top=6pt, bottom=6pt,
  borderline west={2.5pt}{0pt}{accent},
  before skip=10pt, after skip=14pt,
  fontupper=\fontsize{12}{16.5}\selectfont,
  before upper={\wflabel{accentdark}{LA PREGUNTA ECONÓMICA}}}

\newtcolorbox{datosblock}{
  blanker, breakable, colback=tint, boxsep=0pt,
  left=10pt, right=10pt, top=8pt, bottom=8pt,
  before skip=12pt, after skip=14pt,
  fontupper=\small,
  before upper={\wflabel{ink}{DATOS UTILIZADOS}}}

\newtcolorbox{aprendizajesblock}{
  blanker, breakable, boxsep=0pt,
  left=0pt, right=0pt, top=4pt, bottom=4pt,
  before skip=10pt, after skip=12pt,
  before upper={\wflabel{ink}{QUÉ VAS A APRENDER}}}

\newtcolorbox{ideasblock}{
  blanker, breakable, boxsep=0pt,
  left=0pt, right=0pt, top=6pt, bottom=6pt,
  borderline north={1.2pt}{0pt}{ink},
  before skip=14pt, after skip=12pt,
  before upper={\wflabel{ink}{IDEAS CLAVE}}}

% ---------- Leyendas, flotantes y enlaces ----------
\makeatletter
\AtBeginDocument{%
  \captionsetup{font={sf,small,color=ink},labelfont={sf,bf},labelsep=period,
                justification=raggedright,singlelinecheck=false,skip=6pt}%
  \@ifundefined{c@algoritmo}{}{%
    \floatstyle{ruled}\restylefloat{algoritmo}%
    \captionsetup[algoritmo]{font={sf,small},labelfont={sf,bf},labelsep=period,skip=4pt}%
    % Quarto centra el contenido de los flotantes; en un algoritmo estorba.
    \AtBeginEnvironment{algoritmo}{\let\centering\relax}}%
  \hypersetup{colorlinks=true,linkcolor=accentdark,citecolor=accentdark,
              urlcolor=accent,filecolor=accent,linktoc=all}}
\makeatother

% ---------------------------------------------------------------------
% Portada (diseño Warm Folio, versión de septiembre de 2026).
% Un solo motivo: tres densidades sobre una misma base, de la difusa a la
% concentrada, con un filete en la moda. Sin imágenes ni degradados.
% ---------------------------------------------------------------------
\definecolor{densuno}{HTML}{E4C6B4}
\definecolor{densdos}{HTML}{CE8C70}
\definecolor{denstres}{HTML}{B0472A}

\newcommand{\wfportadatitulo}{Macroeconometría\\aplicada con Python}
\newcommand{\wfportadasubtitulo}{De los datos macro a la inferencia estructural y el pronóstico}
\newcommand{\wfportadaautor}{Renato Vassallo}
\newcommand{\wfportadaeyebrow}{LIBRO ABIERTO DE POSGRADO \textperiodcentered\ EDICIÓN 2026}

\renewcommand{\maketitle}{%
  \begin{titlepage}%
  \thispagestyle{empty}%
  \begin{tikzpicture}[remember picture, overlay, x=1in, y=1in]
    % fondo
    \fill[tint] (current page.south west) rectangle (current page.north east);
    % bloque de texto superior
    \node[anchor=north west, xshift=1.05in, yshift=-1.95in, inner sep=0pt,
          font=\sffamily\bfseries\fontsize{9}{11}\selectfont, text=accent]
      at (current page.north west) {\wfportadaeyebrow};
    \draw[ink, line width=0.5pt]
      ([xshift=1.05in, yshift=-2.12in] current page.north west) --
      ([xshift=-1.05in, yshift=-2.12in] current page.north east);
    \node[anchor=north west, xshift=1.02in, yshift=-2.45in, inner sep=0pt, align=left,
          font=\rmfamily\bfseries\fontsize{34}{41}\selectfont, text=ink]
      at (current page.north west) {\wfportadatitulo};
    \node[anchor=north west, xshift=1.05in, yshift=-4.15in, inner sep=0pt, align=left,
          text width=5.0in, font=\rmfamily\fontsize{15}{20}\selectfont, text=muted]
      at (current page.north west) {\wfportadasubtitulo};
    % motivo: tres densidades
    \path (current page.south west) ++(1.05in, 2.75in) coordinate (base);
    \begin{scope}[shift={(base)}]
      \fill[densuno] plot[domain=0:4.9, samples=170, smooth]
        (\x, {1.029*exp(-0.5*((\x-2.058)/1.274)^2)}) -- (4.9,0) -- (0,0) -- cycle;
      \fill[densdos] plot[domain=0:4.9, samples=170, smooth]
        (\x, {1.617*exp(-0.5*((\x-2.450)/0.637)^2)}) -- (4.9,0) -- (0,0) -- cycle;
      \fill[denstres] plot[domain=0:4.9, samples=170, smooth]
        (\x, {2.254*exp(-0.5*((\x-2.695)/0.2695)^2)}) -- (4.9,0) -- (0,0) -- cycle;
      \draw[ink, line width=0.5pt] (2.695,0) -- (2.695,2.377);
      \draw[ink, line width=0.7pt] (0,0) -- (4.9,0);
    \end{scope}
    % pie
    \node[anchor=south west, xshift=1.05in, yshift=1.55in, inner sep=0pt,
          font=\rmfamily\fontsize{15}{18}\selectfont, text=ink]
      at (current page.south west) {\wfportadaautor};
    \node[anchor=south west, xshift=1.05in, yshift=1.25in, inner sep=0pt,
          font=\sffamily\fontsize{8.5}{11}\selectfont, text=muted]
      at (current page.south west) {CC BY-NC-SA 4.0 \textperiodcentered\ compilado el \today};
    \node[anchor=south east, xshift=-1.05in, yshift=1.25in, inner sep=0pt, align=right,
          font=\ttfamily\fontsize{8.5}{13}\selectfont, text=accent]
      at (current page.south east)
      {VAR \textperiodcentered\ BVAR \textperiodcentered\ SVAR\\espacio de estados \textperiodcentered\ ML};
  \end{tikzpicture}
  \end{titlepage}%
}
